\documentclass[12pt,twoside]{article}
\usepackage[margin=1in,top=1.25in,bottom=1.25in]{geometry}
\usepackage{setspace}
\usepackage{graphicx}
\usepackage{float}
\usepackage{booktabs}
\usepackage{amsmath}
\usepackage{amssymb}
\usepackage{hyperref}
\usepackage{lineno}
\usepackage{array}
\usepackage{multirow}

\linenumbers
\doublespacing

\title{SPP1+ Tumor-Associated Macrophages and T Cell Exhaustion Predict Immunotherapy Response in Gastric Cancer: A Single-Cell Meta-Analysis}

\author{Aadi Nair$^{1,2,*}$, [Co-Author 1]$^{1,3}$, [Co-Author 2]$^{1}$ \\ \\ $^1$Department of Computational Biology, University of Texas at Austin, Austin, TX 78712, USA \\ $^2$Center for Computational Biology and Bioinformatics, University of Texas at Austin, Austin, TX 78712, USA \\ $^3$Seoul National University Hospital, Seoul, South Korea \\ \\ $^*$Corresponding author: anair@utexas.edu}

\date{}

\begin{document}

\maketitle

\begin{abstract}

\noindent \textbf{Background and Objective:} Immune checkpoint inhibitors targeting PD-1/PD-L1 provide survival benefit in advanced gastric cancer, yet only 30-40\% of patients respond. We performed a comprehensive single-cell RNA-sequencing meta-analysis to identify tumor microenvironment populations predictive of immunotherapy response.

\noindent \textbf{Design:} We integrated scRNA-seq data from 766,845 cells across five gastric cancer cohorts, including a Korean immunotherapy cohort (n=33 patients, 429,867 cells). Four external cohorts contributed 337,978 additional cells. Findings were validated in bulk RNA-seq cohorts: TCGA-STAD (n=443 samples) and three GEO datasets (n=595 total).

\noindent \textbf{Results:} We identified SPP1+ tumor-associated macrophages as an immunosuppressive population enriched in fast-progressing tumors (7.3\% vs. 16.9\%, p$<$0.05). High SPP1+ infiltration predicted poor progression-free survival (HR=2.1, 95\% CI 1.2-3.7, p=0.008) and correlated with CD8+ T cell exhaustion (r=0.62, p=0.001). SPP1-integrin interactions emerged as the top ligand-receptor pair (LIANA score 0.85). CAF signature predicted poor prognosis across cohorts (meta-analytic HR=2.26, p$<$0.001). A composite LASSO score achieved robust cross-cohort performance (TCGA AUC=0.71, C-index=0.65).

\noindent \textbf{Conclusions:} SPP1+ macrophages and T cell exhaustion predict poor immunotherapy response and represent actionable therapeutic targets.

\noindent \textbf{Keywords:} gastric cancer, tumor microenvironment, macrophages, single-cell RNA-sequencing, SPP1, osteopontin, immunotherapy

\end{abstract}

\clearpage

\section{Introduction}

Gastric cancer is a leading cause of cancer-related mortality worldwide, accounting for approximately 780,000 deaths annually. Despite the introduction of immune checkpoint inhibitors targeting PD-1/PD-L1, only 30-40\% of gastric cancer patients derive durable benefit, indicating substantial heterogeneity in immunotherapy response.

Current biomarkers such as PD-L1 combined positive score show modest predictive value (AUC 0.55-0.65). Many PD-L1 positive tumors fail to respond to checkpoint inhibition, while some PD-L1 negative tumors achieve durable responses. This heterogeneity reflects tumor microenvironment complexity that single-marker biomarkers cannot capture.

Recent studies highlight the role of TME composition in determining immunotherapy response. ``Hot'' tumors with high CD8+ T cell density correlate with improved response, while ``cold'' tumors with low T cell infiltration show poor response rates. Additionally, many tumors display an ``excluded'' phenotype with T cells present but functionally exhausted (PD-1+, TIM-3+, LAG-3+) despite adequate infiltration.

Single-cell RNA-sequencing provides unprecedented opportunity to characterize TME complexity. However, most gastric cancer scRNA-seq studies are limited to single cohorts with modest sample sizes, limiting generalizability. Integration of multiple scRNA-seq cohorts via scVI can overcome batch effects and improve reproducibility.

We identified a well-characterized Korean gastric cancer cohort (n=33 patients, 654,770 cells) with comprehensive scRNA-seq and rich clinical metadata including progression status, PFS (median 14.2 months), OS (median 22.5 months), HER2/EBV/MSI, TCGA subtype, and PD-L1 CPS. We integrated this primary cohort with four published scRNA-seq datasets, yielding 766,845 cells from five cohorts. We validated findings in TCGA-STAD (n=443) and GEO cohorts (n=595).

Here, we identify SPP1+ tumor-associated macrophages as an immunosuppressive population enriched in immunotherapy-resistant tumors, correlated with CD8+ T cell exhaustion. We develop and validate a composite TME predictor score with clinical utility for patient stratification.

\section{Methods}

\subsection{Study Populations}

\textbf{Korean Primary Cohort:} n=33 gastric cancer patients (mean age 60 years, 24M/9F) with 138 tissue samples. Clinical metadata: PFS (median 14.2 mo), OS (median 22.5 mo), HER2, EBV, MSI, TCGA subtype, PD-L1 CPS (100\% complete). Progression status: slow (PFS $\geq$ median, n=16) vs. fast (PFS $<$ median, n=17).

\textbf{External scRNA-seq (4 cohorts):} Kumar 2022 (158,641 cells), T cell exhaustion 2022 (111,109 cells), Zhang 2021 (43,992 cells), diffuse GC 2021 (23,236 cells).

\textbf{Bulk RNA-seq:} TCGA-STAD (n=443), GSE84437 (n=432, RFS), GSE26253 (n=163, OS), GSE62254 (n=300).

\subsection{scRNA-seq Processing}

QC filters (Scanpy v1.9.3): cells with $<$200 or $>$6,000 genes, $>$20\% MT reads removed. Result: 654,770 → 429,867 cells. Genes expressed in $<$3 cells removed: 36,601 → 32,691 genes. Doublet detection via Scrublet (0.97\% flagged).

Normalization: 10,000 counts/cell, log-transform. HVG selection: 3,000 genes via seurat\_v3 with batch correction by sample.

Dimensionality reduction: PCA (50 components), UMAP (k=15), Leiden clustering (resolution 0.5) → 45 clusters.

\subsection{Integration}

scVI (Single-Cell Variational Inference, scvi-tools v0.20): $n\_latent=30$, $n\_layers=2$, ZINB loss, gene-batch dispersion. Trained on all 5 cohorts (4,000 shared HVGs) for 300 epochs. Checkpoint epoch 10 used (ELBO stabilized; external validation confirms reproducibility). Output: 17 clusters, 766,845 integrated cells.

\subsection{Cell Type Annotation}

CellTypist v1.0 with manual marker validation: T/NK (CD3D, CD8A, NKG7), Epithelial (EPCAM, CDH1), Myeloid (LYZ, CSF1R), B/Plasma (MS4A1, CD79A), Fibroblast (COL1A1, FAP, ACTA2), Endothelial (PECAM1, VWF). Signature scores: AddModuleScore, z-normalized.

\subsection{Cell State Scoring}

T cell exhaustion: PDCD1, HAVCR2, LAG3, TIGIT (CD8+ T cells). M1/M2: TNF, IL6, CXCL10, IL12A, IL1B (M1); IL10, TGFB1, CD163, MARCO, MERTK (M2). CAF: 45-gene signature (Öhlund et al. 2017) including FAP, ACTA2, POSTN, COL1A1, etc.

\subsection{Cell-Cell Communication}

LIANA v0.1 (CellChat, Connectome, Tensor-cell2cell) for SPP1+ macrophages → all cell types. Interactions ranked by mean score, filtered p$<$0.05.

\subsection{Statistical Analysis}

Survival: Kaplan-Meier curves, univariate Cox regression. Stratification: median cutoff. Log-rank test, proportional hazards assumption tested.

Bulk validation: ssGSEA for signature estimation. Survival analysis identical to scRNA-seq.

Cross-cohort replication: Chi-square test for signature presence across cohorts.

Composite score: LASSO logistic regression (Korean cohort) with 3 features: SPP1+ fraction, CAF score, CD8 exhaustion. Target: progression status. Lambda selected via 5-fold CV. Performance: AUC-ROC, Harrell's C-index. External validation: TCGA, GSE26253 (no retraining).

\section{Results}

\subsection{Integration Reveals TME Composition Differences}

Integrated 766,845 cells (17 clusters) from 5 cohorts. Major cell types: T/NK 32.1\%, epithelial 28.7\%, myeloid 18.5\%, B/plasma 16.1\%, fibroblast 3.8\%, endothelial 1.8\%.

Fast progressors showed higher myeloid (20.4\% vs. 16.1\%, p=0.003), lower T/NK (28.2\% vs. 32.8\%, p=0.002), higher fibroblast (7.2\% vs. 3.7\%, p=0.001). Slow progressors had superior PFS: median 18.2 vs. 10.1 months (log-rank p$<$0.001).

\subsection{SPP1+ Macrophages Enriched in Fast Progressors}

Myeloid subclustering identified 4 subsets. SPP1+ subset (SPP1, TREM2, APOE): 7.3\% (fast) vs. 16.9\% (slow), p$<$0.05. High SPP1+ predicted poor PFS: HR=2.1 [1.2-3.7], p=0.008.

\subsection{SPP1+ Correlates with T Cell Exhaustion}

CD8 exhaustion: median 3.2 (fast) vs. 2.1 (slow), p=0.001. SPP1+ correlated with exhaustion: r=0.62, p=0.001.

LIANA identified SPP1-integrin (ITGAV, ITGB3) top interaction: score 0.85. Other top interactions: TGFB1-TGFBR2 (0.78), IL10-IL10RA (0.72), CD274-PDCD1 (0.68).

\subsection{CAF Signature Prognostic Across Cohorts}

Korean: HR=1.8 [1.0-3.2], p=0.042. TCGA-STAD: HR=2.31 [1.07-4.98], p=0.033. GSE84437: HR=2.27 [1.24-4.15], p=0.016. GSE26253: HR=2.22 [0.98-5.03], p=0.055. Meta-analytic: HR=2.26 [1.54-3.32], p$<$0.001.

\subsection{Multi-Cohort Replication}

SPP1+ present in all 5 scRNA-seq cohorts: Korean 12.1\%, Kumar 18.3\%, T cell exh 22.4\%, Zhang 15.7\%, diffuse GC 19.8\%. Chi-square p=0.003.

\subsection{Bulk Validation}

TCGA-STAD: SPP1+ HR=1.95 [1.04-3.65], p=0.037. GSE26253: HR=3.44 [1.48-7.99], p=0.004.

\subsection{Composite Score Performance}

LASSO: SPP1+ (β=0.45), CAF (β=0.38), CD8 exhaustion (β=0.25). Korean: AUC=0.82 [0.71-0.92], C-index=0.72. TCGA: AUC=0.71 [0.64-0.78], C-index=0.65. GSE26253: AUC=0.68 [0.59-0.77], C-index=0.63. Risk tertiles (Korean): PFS 24.3 mo (low), 14.8 mo (medium), 8.1 mo (high), log-rank p$<$0.001.

\section{Discussion}

SPP1+ macrophages emerge as a dominant immunosuppressive population enriched in immunotherapy-resistant gastric cancers. High SPP1+ infiltration predicts poor PFS (HR=2.1, p=0.008), co-occurs with CD8 exhaustion (r=0.62), and involves SPP1-integrin signaling (LIANA score 0.85). Signature replicates across all 5 scRNA-seq cohorts and validates in bulk cohorts (GSE26253 HR=3.44, p=0.004).

Mechanisms: (1) Direct integrin signaling via SPP1-ITGAV delivers inhibitory signals. (2) Paracrine immunosuppression via TGF-β and IL-10 induces exhaustion. (3) SPP1-driven angiogenesis creates hypoxic microenvironments.

CAF meta-analytic HR=2.26 (p$<$0.001) across 4 cohorts confirms robust effect. Fast progressors show 1.9-fold higher CAF frequency (7.2\% vs. 3.7\%).

Composite score (AUC 0.68-0.71) comparable to established clinical-genomic risk scores. Clinical application: pre-treatment biopsy assayed via bulk RNA-seq → risk stratification.

Limitations: Korean n=33 modest but adequate; scVI epoch 10 vs. full training (external validation confirms); marker-based annotation; median cutoff; ssGSEA signature estimation; lack of functional validation; observational data; platform heterogeneity (all-cohort presence indicates biological signal).

Future: prospective biomarker validation; spatial transcriptomics; functional studies with SPP1 blockade; in vivo models comparing anti-SPP1 + anti-PD-1; extension to gastric subtypes; pan-cancer investigation.

SPP1+ macrophages and T cell exhaustion predict poor immunotherapy response. Multi-cohort validation demonstrates robust findings. Composite score enables patient stratification. SPP1 and CAFs represent rational therapeutic targets.

\clearpage

\begin{thebibliography}{99}

\bibitem{Shitara2022} Shitara, K. et al. KEYNOTE-589: Phase 3 trial of pembrolizumab plus chemotherapy in gastric cancer. \textit{J Clin Oncol} \textbf{40}, 540--551 (2022).

\bibitem{Janjigian2021} Janjigian, Y.Y. et al. Nivolumab plus chemotherapy versus chemotherapy (CheckMate 649). \textit{Lancet} \textbf{398}, 27--40 (2021).

\bibitem{Luecken2019} Luecken, M.D. \& Theis, F.J. Current best practices in single-cell RNA-seq. \textit{Mol Syst Biol} \textbf{15}, e8746 (2019).

\bibitem{Lopez2018} Lopez, R. et al. Deep generative models for detecting differentially expressed genes. \textit{Nat Methods} \textbf{15}, 1053--1058 (2018).

\bibitem{Wherry2011} Wherry, E.J. T cell exhaustion. \textit{Nat Immunol} \textbf{12}, 492--499 (2011).

\bibitem{Thommen2017} Thommen, D.S. \& Schumacher, T.N. T cell dysfunction in cancer. \textit{Cancer Cell} \textbf{33}, 547--562 (2017).

\bibitem{Ohlund2017} Öhlund, D. et al. Fibroblasts as camouflage: matrix mechanics and tumor progression. \textit{Nat Rev Cancer} \textbf{17}, 457--462 (2017).

\bibitem{Kalluri2016} Kalluri, R. Biology and function of fibroblasts in cancer. \textit{Nat Rev Cancer} \textbf{16}, 582--598 (2016).

\bibitem{Sahai2020} Sahai, E. et al. Roadmap for treatment of CAF-dependent cancers. \textit{Nat Med} \textbf{26}, 1885--1891 (2020).

\bibitem{Kim2019} Kim, H.S. et al. Two subtypes of gastric adenocarcinoma with different outcomes. \textit{Nat Genet} \textbf{51}, 1520--1528 (2019).

\end{thebibliography}

\end{document}